% Ad Familiares 8.4 (ad-familiares-08-04) --- English translation
% Translated by Claude (Anthropic) from the Latin (Perseus).
% Style guide: STYLE.md.

\ciceroLetterOpener{CAELIVS CICERONI S.}

\ciceroSection{1} I envy you. So many things every day are being brought there that would make you stare: first, Messala acquitted; then the same man condemned; C. Marcellus made consul; M. Calidius indicted, after his defeat, by the two Gallii; P. Dolabella made one of the Fifteen. On this last point I do \textit{not} envy you that you missed the prettiest of all spectacles: you did not see Lentulus Crus's face the moment he lost. And with what hope, with what settled conviction he had gone down to the Campus! and with Dolabella himself diffident of his own chances! By Hercules, had not our knights had the sharper eye, he would have all but won, with his rival giving him the ground.

\ciceroSection{2} I do not suppose you were astonished at this: Servaeus, tribune-of-the-plebs designate, condemned, and C. Curio standing for the place in his stead. Curio is putting a great fear into many who do not know him and his easy nature; but as I hope and want, and as he himself behaves, he will prefer the side of the honest men and the Senate; just now, as things stand, he gushes all over with it. The origin and cause of this disposition is that Caesar --- who is in the habit of attaching to himself the friendship of even the lowest sort of men, at whatever expense --- has thoroughly held him in contempt. In which matter it seems to me a most elegantly pretty outcome has fallen out, so much so that the rest have noticed it: Curio, who does nothing by calculation, has appeared to be using calculation and stratagem in evading the schemes which his adversaries had set themselves to lay against his tribunate --- I mean the Laelii and the Antonii and that class of capable men.

\ciceroSection{3} I have sent you this letter after a greater interval because the postponements of the elections were keeping me busier than usual and forcing me to wait day by day for the outcome, so that I might inform you when everything was settled. I waited right up to the Kalends of August. Certain delays have intervened on the praetorian side. As for my own election, what its issue is going to be I do not know. But as far as concerns Hirrus, an unbelievable rumour ran ahead of him at the aedilician elections of the plebs: that bit of fatuity --- the proposal touching the dictatorship, the one we made fun of long ago --- and the publication of the bill suddenly knocked down M. Caelius Vinicianus, and a great shout of laughter pursued him as he fell. From there everyone is now clamouring that Hirrus must not be elected. I hope you will soon hear, both about myself the news you have hoped for, and about that other the news you have scarcely dared to hope for.

\ciceroSection{4} As for the commonwealth, we had already given up looking for anything new; but when a session of the Senate had been held at the temple of Apollo on the 22nd of July, and the matter of Cn. Pompey's pay was being brought forward, mention was made of that legion which Pompey had set down to Caesar's account --- whose force was it now, for as long as Pompey allowed it to be in Gaul? Pompey was forced to say that he would withdraw the legion; not immediately, however, under the goad of those mutterings and outcries of his detractors. From there he was questioned about the succession to C. Caesar; about which --- that is, about the provinces --- it was resolved that Cn. Pompey should return to the city as soon as possible, so that the matter of the succession to the provinces might be transacted in his presence; for Pompey was about to go to Ariminum to the army, and went straight away. I think on the Ides of August the matter will be brought on. Surely either something will be carried through, or there will be a shameful interposition; for in the course of debate Cn. Pompey threw out this remark: that everyone ought to be obedient to the word of the Senate. For my part, however, there is nothing I look forward to so much as Paulus, consul-designate, delivering his opinion first.

\ciceroSection{5} I keep reminding you, over and over, of the Sittian bond (for I want you to understand that the matter is very much my concern); likewise of the panthers --- get hold of those Cibyran fellows, and see to it that they are shipped to me; further, since it is being reported, and is now held for certain, that the king of Alexandria is dead: write me a careful account --- of what you would advise, of how that kingdom now stands, and of who is managing it. The 1st of August.
